\section{机器翻译：从统计方法到神经方法}

\subsection{机器翻译简史}

机器翻译是 NLP 领域最古老的任务之一，也是推动整个领域发展的核心动力。

\begin{figure}[H]
\centering
\includegraphics[width=0.95\textwidth]{slides-images/slide-042.jpg}
\caption{机器翻译的历史：从 1950 年代的规则方法到 1990--2000 年代的统计方法\protect\footnotemark}
\end{figure}
\footnotetext{Slides 第43页。}

\begin{knowledgebox}{机器翻译的起源}
机器翻译的研究可追溯到 1950 年代初——比人工智能作为一个学科诞生还要早。二战期间计算机被用于两件事：（1）计算火炮弹道表；（2）密码破译。冷战期间，人们想到：如果计算机能破解密码，也许也能``破解''语言之间的翻译。大量资金投入这一方向，但在 1950 年代完全失败了——当时人们对语言结构几乎一无所知（Chomsky 的形式语言理论尚未建立），计算机的能力也极其有限。
\end{knowledgebox}

\subsection{统计机器翻译（SMT）}

机器翻译在 1990--2000 年代重新焕发生机，这一时期的主流方法是\textbf{统计机器翻译}（Statistical Machine Translation, SMT）。

\begin{figure}[H]
\centering
\includegraphics[width=0.95\textwidth]{slides-images/slide-046.jpg}
\caption{统计机器翻译系统的特点：极其复杂，需要大量特征工程和人工维护\protect\footnotemark}
\end{figure}
\footnotetext{Slides 第47页。}

SMT 的核心思想是利用贝叶斯公式将翻译问题分解为两个子问题：

$$P(y|x) = \frac{P(x|y) \cdot P(y)}{P(x)} \propto P(x|y) \cdot P(y)$$

\begin{itemize}
    \item $P(x|y)$：\textbf{翻译模型}——词与词之间如何对应翻译（相对简单的词典级模型）
    \item $P(y)$：\textbf{语言模型}——目标语言中什么样的句子是自然的（提供语法和流畅度约束）
\end{itemize}

\begin{warningbox}{SMT 系统的工程复杂度}
Google 的统计机器翻译系统经过十多年、数百名工程师的开发，包含数百万行代码、大量针对特定语言对的手工规则和 hack。每增加一个语言对，都需要重复大量工程工作。系统需要维护短语表、重排序模型、平滑算法等众多独立设计的子组件。
\end{warningbox}

\subsection{翻译中的困难：词序重排与修饰关系}

\begin{figure}[H]
\centering
\includegraphics[width=0.95\textwidth]{slides-images/slide-045.jpg}
\caption{翻译中的词序重排：德语到英语的翻译涉及大量词序变化（Morgen fliege ich $\rightarrow$ Tomorrow I will fly）。中文到英语的翻译更是涉及复杂的修饰关系重组\protect\footnotemark}
\end{figure}
\footnotetext{Slides 第46页。}

Manning 展示了一个经典的中译英示例：``1519年600名西班牙人在墨西哥登陆，去征服几百万人口的阿兹特克帝国，初次交锋他们损兵三分之二。'' Google Translate 在 2009--2015 年间多次翻译这个句子，始终无法正确处理``几百万人口的阿兹特克帝国''这一修饰关系，错误地将``几百万人''理解为征服的主语。

这一例子清楚地展示了 SMT 的根本局限：它无法有效捕捉跨语言的\textbf{句法结构差异}和\textbf{长距离修饰关系}。

\subsection{本章小结}

\begin{itemize}
    \item 机器翻译是 NLP 最古老的任务之一，经历了规则方法（1950s）、统计方法（1990s--2010s）和神经方法（2014--）三个阶段
    \item 统计机器翻译通过贝叶斯分解将问题拆解为翻译模型和语言模型，但工程极其复杂
    \item 不同语言之间的词序重排、修饰关系差异、功能词增减等问题，使翻译远比``逐词替换''复杂
\end{itemize}

%% ========== Seq2Seq ==========

